\documentclass[letterpaper]{article}
\usepackage{aaai26}
\usepackage{times}
\usepackage{helvet}
\usepackage{courier}
\usepackage[hyphens]{url}
\usepackage{graphicx}
\urlstyle{rm}
\def\UrlFont{\rm}
\usepackage{natbib}
\usepackage{caption}
\usepackage{booktabs}
\usepackage{amsmath}
\usepackage{amssymb}
\usepackage{algorithm}
\usepackage{algorithmic}

\frenchspacing
\setlength{\pdfpagewidth}{8.5in}
\setlength{\pdfpageheight}{11in}

\pdfinfo{
/TemplateVersion (2026.1)
}

\setcounter{secnumdepth}{0}

\title{Matched-Gain Diagnostics for On-Policy Distillation: A Minimal Study of Reasoning Stability and Activation-Conditioned Spectra}

\author{
Anonymous Submission
}
\affiliations{
Paper ID: Anonymous
}

\begin{document}

\maketitle

\begin{abstract}
On-policy distillation (OPD) is increasingly used to transfer reasoning behavior from a larger teacher model to a smaller student by supervising the student on states sampled from its own policy.
Yet the empirical question that often matters in practice is not only whether OPD improves an in-domain benchmark, but whether it avoids the out-of-distribution degradation and representation drift that can appear under continued supervised fine-tuning (SFT).
We present a minimal diagnostic protocol for this question.
The protocol compares OPD-like training against continued SFT under matched in-domain gain, and pairs task metrics with activation-conditioned spectral diagnostics computed on frozen prompt and training-target probes.
In a preliminary TRL-based experiment using a Qwen3-1.7B student and Qwen3-4B teacher, an OPD-like run and the nearest SFT control match within 0.010 GSM8K gain.
Under this matched condition, the OPD-like model has lower OOD-lite penalty (0.0100 vs. 0.0182), lower worst OOD-lite drop (0.0100 vs. 0.0150), and substantially lower layer-14 spectral drift averaged across seven target modules (0.0179 vs. 0.0488).
The OPD-like model also has lower drift than the matched SFT control in all seven measured modules.
These results do not establish that OPD improves GSM8K accuracy: evaluation was limited to 200 examples, only one OPD setting was run, and geometry was measured at one layer.
Instead, the contribution is a cautious, reproducible diagnostic framing: under matched task performance, OPD-like training can show a cleaner stability profile than continued SFT, motivating a stricter follow-up with higher-sample evaluation, multiple OPD settings, and multi-layer geometry.
\end{abstract}

\section{Introduction}

Knowledge distillation is a standard way to compress a capable teacher model into a smaller student, but autoregressive distillation creates an immediate mismatch: the student is trained on one set of states and evaluated on states induced by its own generation.
On-policy distillation addresses this mismatch by letting the student generate trajectories and using teacher feedback on the states the student actually visits \citep{agarwal2024gkd}.
Recent OPD work frames this as a mechanism for reducing exposure bias and for progressively aligning student-visited high-probability tokens with teacher-supported tokens \citep{li2026rethinkingopd,song2026surveyopd}.

This paper studies a narrower diagnostic problem.
Suppose a small reasoning model has already received a cold-start distillation phase.
From that same starting point, should one continue with OPD-like training or with ordinary SFT on additional supervised reasoning traces?
A raw accuracy comparison is insufficient.
SFT can move the model toward a narrow trace distribution and may degrade reasoning behavior outside the training style, even when it appears to improve or preserve a nearby benchmark \citep{lobo2025finetuningcot}.
Conversely, OPD can be expensive and unstable when teacher-student overlap is weak \citep{li2026rethinkingopd}.
The useful question is therefore conditional: \emph{under comparable in-domain performance, which training path preserves out-of-distribution behavior and model geometry more cleanly?}

We propose a minimal matched-gain diagnostic for this question.
First, train several continued-SFT controls and at least one OPD-like run from the same starting checkpoint.
Second, match OPD and SFT by gain on an in-domain reasoning benchmark.
Third, evaluate the matched pair using both OOD-lite task metrics and activation-conditioned spectral diagnostics.
The matched-gain step prevents a common confound: a method with lower OOD degradation may simply have learned less.
The spectral step asks whether the task-level pattern is accompanied by a measurable difference in how the trained model changes under probe distributions.

Our preliminary experiment uses Qwen3-1.7B as student, Qwen3-4B as teacher, a NuminaMath-derived prompt pool, TRL-based OPD-like training, and continued SFT controls.
The main matched pair is OPD-like training on 800 rollout prompts and SFT on 256 supervised examples.
Their GSM8K gains differ by 0.010, within the pre-registered valid-match threshold of 0.02.
Under this match, OPD-like training has lower OOD-lite penalty and lower worst OOD-lite drop.
It also has markedly lower activation-conditioned spectral drift from the starting checkpoint at layer 14, both on average and in every measured attention and MLP projection module.

The result should be read carefully.
This is not a full OPD replication, not a claim of broad OOD generalization, and not evidence that OPD improves GSM8K accuracy.
The preliminary run used an evaluation limit of 200 examples, only one OPD hyperparameter setting, and geometry from a single layer.
The value of the study is that it closes a minimal engineering loop and identifies a sharper next experiment: test whether the same stability signal survives lower evaluation noise, additional OPD settings or seeds, a less saturated starting point, and broader geometry.

The contributions are:
\begin{itemize}
    \item A matched-gain diagnostic protocol for comparing OPD-like training and continued SFT under comparable in-domain reasoning performance.
    \item An activation-conditioned spectral analysis that links OOD-lite degradation to drift from the starting checkpoint.
    \item A preliminary TRL-based result showing that OPD-like training can preserve a cleaner stability profile than continued SFT under a valid GSM8K match.
    \item A concrete follow-up gate for moving from a stability claim to a possible improvement-plus-stability claim.
\end{itemize}

\section{Related Work}

\paragraph{On-policy distillation.}
Generalized Knowledge Distillation trains the student on self-generated output sequences and uses teacher feedback on those sequences, directly targeting the distribution mismatch between teacher-generated training outputs and student-generated inference outputs \citep{agarwal2024gkd}.
More recent OPD surveys and mechanism studies organize this family around student-sampled trajectories, feedback granularity, and stability conditions \citep{song2026surveyopd,li2026rethinkingopd}.
The most relevant mechanism claim is that successful OPD aligns high-probability teacher-supported tokens at student-visited states, while poor teacher-student overlap can cause failure.
This motivates the cold-start plus OPD-like setup used here: the student should be close enough to the teacher-supported region that on-policy feedback is meaningful.

\paragraph{Fine-tuning and reasoning degradation.}
SFT is efficient and simple, but its effect on reasoning can be mixed.
Recent work on fine-tuning and chain-of-thought reasoning reports that task-specific fine-tuning can reduce reasoning performance and affect faithfulness, especially for smaller models \citep{lobo2025finetuningcot}.
This does not imply that SFT is always harmful.
It does imply that a continued-SFT baseline should be evaluated for degradation, not treated as a neutral amount of additional training.
Our diagnostic uses several SFT data sizes precisely because the amount of supervised trace fitting can change both in-domain and OOD behavior.

\paragraph{Spectral views of adaptation.}
Spectral analysis has become a useful lens for understanding how fine-tuning changes models.
For example, \citet{shuttleworth2025lora} show that LoRA and full fine-tuning can produce different singular-vector structures, and connect those structures to forgetting.
Our analysis is different: it does not inspect raw weight matrices alone.
Instead, it computes singular spectra after conditioning linear weights by activation statistics collected from probe text.
This makes the diagnostic sensitive to the distribution under which a module is being used.
The broader motivation is the same: if two adaptation procedures obtain similar task performance but move the model through different spectral paths, those paths may explain differences in retention and generalization.

\paragraph{Representation-level OPD.}
OPRD proposes to move on-policy distillation from output-space token distributions into hidden-state alignment \citep{yang2026oprd}.
Our experiment does not implement representation-level supervision.
However, the result points in a compatible direction: if OPD-like output training is associated with lower activation-conditioned drift, then representation-level diagnostics and future representation-level objectives are natural extensions.

\section{Problem Setup}

Let $\theta_0$ denote a cold-start student checkpoint obtained by distilling from a teacher on a small prompt set.
From $\theta_0$, we compare two training families.
The first is OPD-like training, where the current student samples completions and receives token-level teacher supervision on those sampled trajectories.
The second is continued SFT, where the student is trained on fixed supervised reasoning traces.

Let $m(\theta)$ be an in-domain reasoning metric, here GSM8K accuracy, and let
\begin{equation}
g(\theta) = m(\theta) - m(\theta_0)
\end{equation}
be the gain over the common starting checkpoint.
For each OPD-like model $\theta_o$, we select the nearest SFT model $\theta_s$ by absolute gain gap:
\begin{equation}
\Delta_g(\theta_o,\theta_s)=|g(\theta_o)-g(\theta_s)|.
\end{equation}
A pair is treated as a valid match when $\Delta_g \leq 0.02$.

The central stability question is then:
\begin{quote}
Among models with comparable in-domain reasoning gain, does OPD-like training produce lower OOD-lite degradation and lower activation-conditioned spectral drift than continued SFT?
\end{quote}

This framing deliberately separates two claims.
The first is a stability claim: OPD preserves behavior and geometry better than a matched SFT control.
The second is a stronger improvement claim: OPD improves the in-domain benchmark while preserving stability.
The present experiment can only support the first claim.

\section{Method}

\subsection{Training Arms}

The experiment uses a single prompt pool split into training, held-out, and probe subsets.
A Qwen3-1.7B student is first cold-started from a Qwen3-4B teacher using 512 prompts.
From this $\theta_0$, we train one OPD-like arm and four continued-SFT controls.
The OPD-like arm uses $\lambda=1.0$ and consumes 800 rollout prompts.
The SFT controls use 256, 512, 1024, and 2048 supervised examples.
All arms use the same LoRA target modules and start from the same $\theta_0$.

\subsection{Task Metrics}

The in-domain matching metric is GSM8K.
Because the preliminary run used an evaluation limit of 200 examples, differences smaller than roughly 0.02 should not be interpreted as reliable improvements.
OOD-lite performance is summarized by an average OOD score, a penalty $p2$, and a worst-drop metric.
The exact task mix follows the local Cycle 03 protocol; in this earliest paper draft, the important point is not the absolute OOD score but the matched comparison between OPD-like and SFT models.

\subsection{Activation-Conditioned Spectral Diagnostics}

For each selected layer and target linear module, we collect activation statistics on probe text and form a whitened or activation-conditioned weight matrix.
Let $W$ be the module weight matrix and $L_{\mathcal{D}}$ be the Cholesky factor derived from the activation Gram matrix on probe distribution $\mathcal{D}$.
We analyze the singular spectrum of
\begin{equation}
A^{\mathcal{D}} = W L_{\mathcal{D}}.
\end{equation}
This spectrum captures how strongly the module acts along directions that are active under the probe distribution.

Given singular values $\sigma=(\sigma_1,\ldots,\sigma_n)$, the effective rank is
\begin{equation}
\operatorname{erank}(\sigma)=\exp\left(-\sum_i p_i \log p_i\right), \quad
p_i=\frac{\sigma_i}{\sum_j \sigma_j}.
\end{equation}
The drift from the starting checkpoint is the root mean squared log-spectrum distance:
\begin{equation}
\operatorname{drift}(\sigma,\sigma^0)
=\sqrt{\frac{1}{L}\sum_{i=1}^L
\left(\log(\sigma_i+\epsilon)-\log(\sigma_i^0+\epsilon)\right)^2}.
\end{equation}
We also track the spectral gap and the signed log-spectrum level gap between a frozen reference probe $X$ and each model's training-target probe $S$.

\begin{algorithm}[t]
\caption{Matched-Gain Stability Diagnostic}
\begin{algorithmic}[1]
\STATE Train common cold-start checkpoint $\theta_0$.
\STATE Train OPD-like arms $\{\theta_o\}$ from $\theta_0$.
\STATE Train SFT controls $\{\theta_s\}$ from $\theta_0$.
\STATE Evaluate GSM8K and compute $g(\theta)=m(\theta)-m(\theta_0)$.
\FOR{each OPD-like arm $\theta_o$}
    \STATE Select $\theta_s^\star = \arg\min_{\theta_s}|g(\theta_o)-g(\theta_s)|$.
    \STATE Mark valid match if $|g(\theta_o)-g(\theta_s^\star)|\leq 0.02$.
    \STATE Compare OOD-lite degradation under the valid match.
    \STATE Compare activation-conditioned spectral drift under the valid match.
\ENDFOR
\end{algorithmic}
\end{algorithm}

\section{Preliminary Experiment}

\subsection{Configuration}

Table~\ref{tab:config} summarizes the current minimal run.
The design is intentionally small: it verifies whether the end-to-end OPD-like training, evaluation, matching, and geometry pipeline can close before larger experiments are run.

\begin{table}[t]
\centering
\small
\begin{tabular}{ll}
\toprule
Item & Value \\
\midrule
Student & Qwen3-1.7B \\
Teacher & Qwen3-4B \\
Cold start & 512 prompts, OPD-like distillation \\
OPD-like arm & $\lambda=1.0$, 800 rollout prompts \\
SFT controls & 256, 512, 1024, 2048 examples \\
Main eval & GSM8K, limit 200 \\
Geometry & Layer 14, 7 target modules \\
Probe length & 512 tokens \\
LoRA & rank 16, alpha 32, dropout 0.05 \\
\bottomrule
\end{tabular}
\caption{Minimal Cycle 03 configuration.}
\label{tab:config}
\end{table}

\subsection{Task Trajectory}

Table~\ref{tab:trajectory} reports the task trajectory.
The most important observation is that larger SFT controls degrade GSM8K under this setup, while the OPD-like arm approximately preserves the starting checkpoint's GSM8K score.
Because the evaluation limit is only 200, the OPD-like score should be interpreted as maintenance, not improvement.

\begin{table}[t]
\centering
\small
\begin{tabular}{lrrrr}
\toprule
Model & Size & GSM8K & Gain & OOD $p2$ \\
\midrule
SFT-1024 & 1024 & 0.295 & -0.140 & 0.0261 \\
SFT-2048 & 2048 & 0.310 & -0.125 & 0.0279 \\
SFT-512 & 512 & 0.380 & -0.055 & 0.0124 \\
OPD-$\lambda1$ & 800 & 0.430 & -0.005 & 0.0100 \\
$\theta_0$ & 512 & 0.435 & 0.000 & 0.0000 \\
SFT-256 & 256 & 0.440 & 0.005 & 0.0182 \\
\bottomrule
\end{tabular}
\caption{Evaluation trajectory. Size denotes consumed rollout prompts for OPD-like training and supervised examples for SFT.}
\label{tab:trajectory}
\end{table}

\subsection{Matched-Gain Comparison}

The nearest SFT match for OPD-$\lambda1$ is SFT-256.
The GSM8K gain gap is
\[
|-0.005 - 0.005| = 0.010,
\]
which satisfies the valid-match threshold.
Table~\ref{tab:matched} shows the matched OOD-lite comparison.
Under matched GSM8K gain, the OPD-like model has lower OOD-lite penalty and lower worst OOD-lite drop.

\begin{table}[t]
\centering
\small
\begin{tabular}{lrrr}
\toprule
Metric & OPD & SFT-256 & OPD--SFT \\
\midrule
GSM8K gain & -0.005 & 0.005 & -0.010 \\
OOD-lite penalty $p2$ & 0.0100 & 0.0182 & -0.0082 \\
Worst OOD-lite drop & 0.0100 & 0.0150 & -0.0050 \\
\bottomrule
\end{tabular}
\caption{Matched-gain comparison. Negative OPD--SFT values are favorable for degradation metrics.}
\label{tab:matched}
\end{table}

\subsection{Spectral Stability}

Table~\ref{tab:geometry_summary} summarizes layer-14 geometry.
The OPD-like model has a spectral drift of 0.0179, compared with 0.0488 for SFT-256.
The SFT controls show increasing drift as the amount of supervised training data increases.
The OPD-like model also keeps the $X$--$S$ spectrum level gap closer to zero than the SFT controls, suggesting that the reference generation distribution and training-target distribution remain more aligned under the OPD-like update.

\begin{table}[t]
\centering
\small
\begin{tabular}{lrrr}
\toprule
Model & Eff. rank & Drift & $X$--$S$ gap \\
\midrule
$\theta_0$ & 897.29 & 0.0000 & -0.0231 \\
OPD-$\lambda1$ & 892.66 & 0.0179 & -0.0185 \\
SFT-256 & 925.28 & 0.0488 & -0.0467 \\
SFT-512 & 927.27 & 0.0527 & -0.0467 \\
SFT-1024 & 930.36 & 0.0598 & -0.0479 \\
SFT-2048 & 930.84 & 0.0643 & -0.0492 \\
\bottomrule
\end{tabular}
\caption{Layer-14 activation-conditioned spectral summary averaged across seven target modules.}
\label{tab:geometry_summary}
\end{table}

The module-level result is stronger than the average alone.
As shown in Table~\ref{tab:module_drift}, OPD-like training has lower drift than SFT-256 in every measured layer-14 module.
The largest absolute differences appear in the attention output projection and MLP gate/up projections.
This pattern is consistent with the hypothesis that continued SFT pushes the model toward a narrower answer-style manifold, while OPD-like training performs a milder distribution-aligned update.

\begin{table}[t]
\centering
\small
\begin{tabular}{lrrr}
\toprule
Module & OPD & SFT-256 & Diff. \\
\midrule
attn.q & 0.0165 & 0.0391 & -0.0226 \\
attn.k & 0.0137 & 0.0355 & -0.0218 \\
attn.v & 0.0137 & 0.0357 & -0.0220 \\
attn.o & 0.0328 & 0.0792 & -0.0464 \\
mlp.gate & 0.0162 & 0.0597 & -0.0435 \\
mlp.up & 0.0158 & 0.0597 & -0.0439 \\
mlp.down & 0.0167 & 0.0324 & -0.0157 \\
\bottomrule
\end{tabular}
\caption{Layer-14 module drift from $\theta_0$. Negative differences mean lower OPD drift.}
\label{tab:module_drift}
\end{table}

\section{Discussion}

The preliminary result supports a conservative version of the K1 hypothesis:
under matched in-domain reasoning gain, OPD-like training can produce lower OOD-lite degradation and lower activation-conditioned spectral drift than continued SFT.
The matched-gain requirement is important.
Without it, lower drift could be dismissed as a sign that OPD trained less.
Here, OPD-$\lambda1$ and SFT-256 are close in GSM8K gain, but their OOD-lite and spectral profiles differ.

The result also clarifies what is not yet known.
First, OPD-like training did not improve GSM8K over $\theta_0$ in a reliable way.
Second, the starting checkpoint may be saturated: if $\theta_0$ already reaches the local GSM8K ceiling under this limited evaluation, an improvement claim is difficult to detect.
Third, the geometry evidence is only layer-14 evidence.
The module-level consistency is encouraging, but a mechanism claim needs multi-layer evidence and ideally singular-vector or principal-angle artifacts.

These constraints motivate the next experiment rather than weakening the usefulness of the current one.
A stricter follow-up should use a less saturated $\theta_0$, increase the GSM8K evaluation sample count or remove the limit, include at least one additional OPD setting or seed, and measure geometry across early, middle, and late layers.
If the stability pattern persists, the work can support a stronger claim that OPD-like state matching controls destructive drift relative to continued SFT.
If the same setup also produces a reliable GSM8K gain, the claim can move from stability to improvement plus stability.

\section{Limitations}

This earliest draft has several limitations that must be resolved before a submission-quality AAAI paper.
The experiment has only one OPD-like run.
The GSM8K evaluation limit of 200 makes small differences unreliable.
The OOD-lite suite is useful as a local degradation probe but does not establish broad OOD generalization.
The spectral analysis is limited to layer 14 and does not yet include stored singular-vector sketches or principal-angle analysis.
The implementation is TRL OPD-like rather than a full reproduction of a standard OPD system such as a verl implementation.
Finally, the current evidence supports stability under matched performance, not a general performance improvement claim.

\section{Conclusion}

We presented a minimal matched-gain diagnostic for comparing OPD-like training with continued SFT.
In a preliminary Qwen3-1.7B/Qwen3-4B experiment, OPD-like training matched the nearest SFT control on GSM8K gain while showing lower OOD-lite degradation and substantially lower activation-conditioned spectral drift.
The result is best viewed as a feasibility and stability signal, not as a final OPD performance claim.
The next research gate is clear: test the same matched-gain stability pattern under lower evaluation noise, multiple OPD settings, a less saturated starting checkpoint, and multi-layer geometry.

\bibliography{references}
\bibliographystyle{aaai26}

\end{document}

